% Source: https://artofproblemsolving.com/community/c3h1535032_trapezium_medians_parallel_proof
% Stable identity: aops-topic-1535032-q01
Draw the diagonal $AC$ and let $K$ be its midpoint:

\begin{asy}
usepackage("amsmath");
size(6cm);
defaultpen(fontsize(10pt));
pair A=(0,0),B=(7,0),C=(5.2,3.2),D=(1.3,3.2);
pair M=(A+D)/2,Nn=(B+C)/2;
draw(A--B--C--D--cycle,linewidth(1));
draw(M--Nn,red+linewidth(1));
draw(A--C,dashed+gray(0.4));
pair K=(A+C)/2; dot(K,blue); label("$K$",K,dir(-80),blue);
dot(A);dot(B);dot(C);dot(D);dot(M,red);dot(Nn,red);
label("$A$",A,SW);label("$B$",B,SE);label("$C$",C,NE);label("$D$",D,NW);
label("$M$",M,W,red);label("$N$",Nn,E,red);
label("\mbox{}",(min(currentpicture,user=true)+max(currentpicture,user=true))/2,black); // keeps the text after the figure black when LaTeX typesets the labels inline
\end{asy}

\textbf{Triangle $ACD$.} $M$ and $K$ are the midpoints of $AD$ and $AC$, so by the midpoint theorem (triangle $AMK$ is similar to triangle $ADC$ with ratio $\dfrac{1}{2}$, by SAS similarity, since they share angle $A$ and $\dfrac{AM}{AD}=\dfrac{AK}{AC}=\dfrac{1}{2}$) we get $MK\parallel DC$ and $MK=\dfrac{1}{2}DC$.

\textbf{Triangle $CAB$.} In the same way $K$ and $N$ are the midpoints of $CA$ and $CB$, so $KN\parallel AB$ and $KN=\dfrac{1}{2}AB$.

\textbf{Conclusion.} Since $AB\parallel DC$, the lines $MK$ and $KN$ are both parallel to the bases and pass through the same point $K$, so they are the same line: $M$, $K$, $N$ are collinear and $MN\parallel AB\parallel DC$. Moreover $K$ lies between $M$ and $N$ ($M$ and $N$ are on opposite sides of the line $AC$, being midpoints of the sides $AD$ and $BC$, which lie on opposite sides of the diagonal), so
\[
MN=MK+KN=\dfrac{DC+AB}{2}.\qquad\blacksquare
\]
